
% ============================================================
% SECTION 01
% THE ONTOLOGICAL PROBLEM OF INVARIANCE
% Chapter 02: Invariance
% Book I: Ontology of Viability
% ============================================================

\section{The Ontological Problem of Invariance}
\label{sec:ontological-problem-of-invariance}

% ------------------------------------------------------------
% 1.1. THE QUESTION OF INVARIANCE
% ------------------------------------------------------------

The problem of invariance emerges when an existent is
considered in relation to a domain of difference or change.
An existent may be subject to different conditions,
transformations, or alterations, while something about it
remains unchanged in a particular respect. This possibility
raises a question that cannot be adequately addressed by
examining identity alone: what does it mean for something
to remain invariant in relation to change?

The question is ontological because it concerns the
determinations of an existent and the manner in which
these determinations stand in relation to possible changes.
It is not merely a question about the ability of an observer
to recognize similarities between different states, nor is
it limited to the identification of recurring patterns.
The epistemic recognition of invariance is one dimension
of the problem, but it does not exhaust its ontological
significance.

The problem must therefore be approached without assuming
that everything which remains unchanged is necessarily
the identity of the existent, or that everything which
changes necessarily constitutes a transformation of its
identity. Both assumptions would prematurely collapse
distinct ontological questions into one another.

The first task is to clarify the meaning of invariance
itself and to identify the conditions under which an
invariance claim can be meaningfully formulated.

% ------------------------------------------------------------
% 1.2. THE RELATIONAL CHARACTER OF INVARIANCE
% ------------------------------------------------------------

For the purposes of this first edition, invariance is
provisionally understood as a relation between an existent,
or a determinate aspect of an existent, and a specified
domain of changes or transformations, such that the
existent or aspect remains unchanged in a particular
respect under those changes or transformations.

This definition is relational because invariance cannot
be adequately specified without identifying the domain
in relation to which something is said to remain unchanged.
A statement that something is invariant without specifying
what it is invariant in relation to remains incomplete.

The relational formulation requires three analytically
distinguishable elements:

\begin{enumerate}
    \item \textbf{The subject of invariance:} the existent
    or determinate aspect whose invariance is being
    considered.

    \item \textbf{The domain of change:} the specified
    transformations, alterations, or variations in relation
    to which invariance is examined.

    \item \textbf{The respect of invariance:} the particular
    determination, feature, or relation that remains
    unchanged within that domain.
\end{enumerate}

These elements define the structure of an invariance claim.
They do not establish that the elements themselves are
independent ontological entities. Their distinction is
analytical and serves to prevent conceptual ambiguity.

Consider, for example, an existent whose spatial position
changes while a particular structural relation among its
components remains unchanged. In this case, the position
does not remain invariant, but the specified structural
relation may do so. The invariance claim is therefore
restricted to a particular determination and a particular
domain of change.

This example illustrates why invariance should not be
understood as the absence of all alteration. It is possible
for some determinations to change while others remain
unchanged in relation to the same transformation.

The example does not, however, establish that structural
relations are the ultimate ground of identity or that
structural invariance is universal. It only demonstrates
the need to specify the respect in which invariance is
being asserted.

% ------------------------------------------------------------
% 1.3. INVARIANCE AND THE PROBLEM OF CHANGE
% ------------------------------------------------------------

The concept of invariance becomes meaningful through its
relation to a specified domain of change. This does not
mean that change is ontologically prior to invariance,
nor that invariance exists only as a consequence of an
observer's comparison between different states.

Rather, the relational definition identifies the terms
required to formulate a claim of invariance. The domain
of change establishes the respect in which the question
of remaining unchanged is being considered.

A distinction must therefore be maintained between
the occurrence of change and the determination of what
remains unchanged in relation to it.

Change may affect different aspects of an existent in
different ways. Some aspects may undergo alteration,
others may remain unchanged within the specified domain,
and the status of still others may remain undetermined.
It would be unjustified to assume that every change
affects every determination of an existent equally.

Likewise, it would be unjustified to infer from the
persistence of one determination that the existent as
a whole remains unchanged.

The problem of invariance consequently requires a
differentiated examination of the existent rather than
a general opposition between change and permanence.

% ------------------------------------------------------------
% 1.4. INVARIANCE IS NOT ABSOLUTE IMMUTABILITY
% ------------------------------------------------------------

Invariance must be distinguished from immutability.
Immutability, in its strongest sense, indicates the
absence of change in the subject under consideration.
Invariance, as provisionally defined here, does not
require such an absence.

An existent may undergo changes in certain respects
while remaining invariant in another respect. The
invariance of a particular determination therefore
does not establish the immutability of the whole existent.

Conversely, the absence of an observed change does not
by itself establish invariance. Such an absence may
result from the limited scope of observation, the
restricted duration of examination, or the fact that
the relevant domain of transformations has not been
specified.

The distinction is important because the concept of
invariance must not be transformed into an unsupported
claim about absolute permanence.

At this stage, the analysis does not determine whether
absolute invariance is possible, whether it has a
specific ontological status, or whether it applies
to any existent in an unrestricted sense. These
questions remain open.

% ------------------------------------------------------------
% 1.5. INVARIANCE AND THE INSEPARABILITY OF BEING
% AND IDENTITY
% ------------------------------------------------------------

The relational definition of invariance must remain
consistent with the foundational commitment established
in the preceding chapter: being and identity are
inseparable.

The introduction of invariance does not imply that
an existent first possesses an identity and subsequently
acquires invariant features as additional ontological
properties. Nor does it imply that invariance constitutes
a prior foundation from which identity emerges.

An existent is considered here as already identity-bearing.
The analysis of invariance concerns a further question:
how particular determinations of that existent stand
in relation to specified changes.

This distinction prevents invariance from being
mistaken for a substitute for identity. The fact that
a determination remains unchanged under a particular
transformation does not, by itself, explain why the
existent is the existent that it is.

Similarly, the fact that an existent undergoes change
does not, by itself, establish that its identity has
ceased to be.

The relation between identity and invariance therefore
requires further examination. The present section
establishes only that the two concepts address distinct
but connected problems.

% ------------------------------------------------------------
% 1.6. THE EPISTEMIC AND ONTOLOGICAL DIMENSIONS
% OF THE PROBLEM
% ------------------------------------------------------------

Claims about invariance are often formulated through
comparison. An observer examines different conditions
or states and identifies something that appears to
remain unchanged. This comparative activity belongs
to the epistemic dimension of the problem.

However, the epistemic act of identifying an invariant
must be distinguished from the ontological status of
what is identified.

The recognition of an invariant is conditioned by
the knower's position, available distinctions, methods
of examination, and limits of knowledge. A feature
may remain unidentified because the knower lacks
the conceptual or empirical resources necessary to
recognize it.

The absence of recognition therefore cannot be treated
as proof of the absence of invariance.

At the same time, an epistemic claim of invariance
must remain open to examination and correction.
What appears invariant within one limited domain
may prove variable when a broader or different domain
of transformations is considered.

Knowledge and non-knowledge are consequently intertwined
in the investigation of invariance. The knower necessarily
occupies a situated position from which some determinations
can be recognized while others remain unknown.

This limitation does not reduce being to knowledge.
It requires a distinction between what is ontologically
the case and what can be established by a particular
act of knowing.

% ------------------------------------------------------------
% 1.7. THE SCOPE OF THE INITIAL DEFINITION
% ------------------------------------------------------------

The relational definition adopted in this chapter
is deliberately limited.

It establishes that a meaningful claim of invariance
must specify:

\begin{itemize}
    \item the existent or determination under consideration;
    \item the domain of changes or transformations;
    \item the respect in which something remains unchanged.
\end{itemize}

It does not establish:

\begin{itemize}
    \item that invariance is the ultimate ground of identity;
    \item that every existent necessarily possesses an
    invariant feature;
    \item that every invariant feature is essential;
    \item that invariance is equivalent to permanence
    or immutability;
    \item that the ground of invariance must be a substance,
    essence, structure, or higher ontological determination;
    \item that invariance has the same form in every
    ontological domain.
\end{itemize}

These restrictions are not deficiencies in the definition.
They are methodological safeguards against introducing
metaphysical commitments before the relevant arguments
have been developed.

The possibility that invariance may possess a deeper
ontological significance beyond its relational
formulation remains open for subsequent investigation
and for future revisions of the philosophical system.

% ------------------------------------------------------------
% 1.8. THE ONTOLOGICAL SIGNIFICANCE OF THE PROBLEM
% ------------------------------------------------------------

The problem of invariance is significant because it
makes possible a more precise examination of the
relationship between an existent and the changes
in relation to which its determinations are considered.

Without the concept of invariance, the analysis may
remain trapped between two inadequate alternatives:
either the existent is assumed to remain entirely
unchanged, or every alteration is treated as if it
necessarily affects the existent in its entirety.

The relational approach makes room for a more
differentiated account. It permits the examination
of what changes, what remains unchanged, and what
cannot yet be determined within a specified domain.

This does not resolve the problem of persistence,
the origin of identity, or the ontological foundation
of continuity. It provides a conceptual framework
within which these questions can be distinguished
rather than prematurely answered.

In this respect, invariance contributes to the
development of the ontology of viability by refining
the analysis of how identity-bearing existents may
stand in relation to difference and transformation.

Its significance lies not in providing an immediate
explanation of persistence, but in clarifying the
conditions under which claims about remaining
unchanged can be formulated and critically examined.

% ------------------------------------------------------------
% 1.9. TRANSITION TO THE FOLLOWING SECTIONS
% ------------------------------------------------------------

Having established the ontological problem and the
provisional relational definition of invariance,
the analysis must now examine its relation to the
concepts developed in the preceding chapter.

The first and most immediate comparison concerns
identity. Since identity and being have already been
established as inseparable, it is necessary to
determine how invariance can be discussed without
being confused with identity or treated as its
ontological foundation.

The following section therefore examines the
relationship between invariance and identity,
identifying their conceptual differences and the
limits of what can be inferred from their relation.

% ============================================================
% END OF SECTION 01
% ============================================================
